\section{Realizing the chosen degrees in the actual family}
\label{sec:host-weight-join}

\status{We prove the degree-realization part of the actual
\texttt{HostWW--KitS--ChosenWall} route. Arithmetic hosts, the point
constructor for every admissible degree pair, and its numerical records
remain supplied inputs. The stronger property \texttt{CoverAll} is not
needed on this route.}

Fix the scale $s$, ranks $R$, host $H$, constants $C,c_D,c_h$ with
\texttt{Std} $C$ and $C.\texttt{GK.Nonneg}$, and an
admissible word code. A supplied constructed point $P$ includes its actual
ties and extra data. Assume \texttt{Val}, \texttt{Pt.From} and
\texttt{CanonInst}, as in the source. Write $T=H.\texttt{TT}$ and
$\Delta=\lceil x^{c_D}\rceil_{\mathbb N}$,
$h_0=2\lceil x^{c_h}\rceil_{\mathbb N}$, and $w''_n=b'/(xn)$.
For integer functions $a,b$ on $R.A$, let
\[
 \mathcal E_H(a,b)=
 \sum_{p\in\mathcal C_H}\frac{\log p}{p}
 \left(2\mathbf1_{\{\Sigma_{H,p}(a,b)>x^{-\epsilon}\}}
       +4\Sigma_{H,p}(a,b)\right).                 \tag{HW.1}
\]
Here $\mathcal C_H$ and $\Sigma_{H,p}$ are the literal host candidate
set and canonical-deficit sum. The source's \texttt{WOK} is
$\mathcal E_H(a,b)\le x^{-9}$; \texttt{WReq} is the point version.

\paragraph{HW1: exact host--point transport.}
The value map on $R.A$ is injective. Thus \texttt{Pt.From} identifies
the selected sets, Frobenius steps, $q,d$, and prime classifications,
not just their cardinalities. The actual numerical bases agree because
\[
 \mathrm{Sc}=4k\Delta,\qquad \eta=mh_0.             \tag{HW.2}
\]
The two definitions of the universal admissible-deficit set
$\mathcal K(p,n',T',a',b')$ are identical after these substitutions.
The point's extra record gives
$\kappa_I=\inf\mathcal K$; separately, \texttt{CanonInst} gives the
canonical orbit selector $\mathrm{bs}=\mathrm{bs}_H$.
Consequently the representatives, local deficits, and every summand in
(HW.1) agree. No choice of a minimizing certificate is involved.

If GW5 supplies \texttt{WeightVectorExists} at $P$, transfer its chosen
functions along this identification. The box has value zero outside
$T$. Its integer endpoints are exactly the inequalities
\[
 a\le a_t\le a+\frac{a}{kx^{13}},\qquad
 b'\le b_t\le b'+\frac{2n\delta x^\epsilon\Lambda b'}k
 \quad(t\in T).                                  \tag{HW.3}
\]
Hence it supplies the host's \texttt{WinOK} and \texttt{WOK}.
Conversely, if these host functions are the degrees in
\texttt{Pt.From}, their zero extensions are the point's actual degrees,
and host \texttt{WOK} gives actual \texttt{WReq} with the same constants.

\paragraph{HW2: choose the pair, then reconstruct every position.}
Assume the actual point constructor is available for \emph{each}
admissible pair and every $u_0<n$, uniformly beyond its supplied
threshold. First use any admissible pair to build temporary points and
kits. With the supplied \texttt{NumAt} records, the family constructor
produces a temporary family to which GW5 applies. Its window membership
gives an index at every sufficiently large allowed scale. HW1 transfers
the selected pair from that point back to its original host. This is the
degree-selection assertion of \texttt{HostWW}.

Keep this new pair fixed and invoke the point constructor again, for
each $u_0\in\{0,\ldots,n-1\}$. These points form an actual kit.
For completeness, its less immediate equality fields can be checked
directly. Their canonical $\kappa_I$ agree because $k,s,\mathrm{Sc}$
determine the same admissible-deficit sets; for every $p\in\mathbb N$
their exponents are the same
$\lceil C_e f k^2\log_p a\rceil_{\mathbb N}$, with the source's
value zero at $p=0,1$. At a lifted level $1\le j\le m$,
the filling equation and the common Hodge profile give
\[
 \operatorname{off}(j,(0,t))=
 \begin{cases}\lfloor b'/x\rfloor,&t\in T,\\
               \lceil w''_n\rceil,&t\notin T
 \end{cases}
 +(j-1)h_0+\frac{h_0}{2}.                         \tag{HW.4}
\]
Outside the lifted levels \texttt{CanonInst} prescribes $b'/x$ or
$w''_n$ directly. The remaining host equalities come from
\texttt{Pt.From}; the common code identifies the word labels. Indexing
by $u_0$ gives \texttt{Kit.cover}, while the fixed pair gives degree
equality at every common value label. HW1 and (HW.3) give
\texttt{WallPt} at \emph{every} reconstructed point.

All constants and construction thresholds are fixed before the scale;
the new pair may depend on that scale and host. The points for all
$u_0$ are chosen only after the pair. Thus the argument does not require
an already fixed family to contain every admissible degree pair.

\paragraph{HW3: actual family transport and an existing carrier.}
Construct \texttt{famOfKits} from these kits and their supplied
\texttt{NumAt} records. It copies each point's $D,L,G_x$, ties and extra
data; its common fields are the kit equalities just proved. Therefore
\texttt{WallPt} holds at every family index. The identity
$2\mathbf1_E=(2\text{ if }E\text{ else }0)$ converts its loss to the
literal \texttt{ChW\_at}; beyond the family threshold this proves
\texttt{ChosenWall}.

For an allowed $x$, choose an existing index $i_0$ using
\texttt{Ix\_win}. Transfer its actual degrees to the common integer
labels $T_n(x)$, assigning zero elsewhere. The point windows and the
family's $b_T\le b_{\rm top}$ give the required admissibility. The
index $i_0$ itself carries this pair, and every other carrier satisfies
\texttt{ChW\_at} because every family member does. This is exactly
\texttt{ChosenWS}. Empty index sets outside the scale window cause no
problem. Arithmetic host existence, the biased top-set
\texttt{Frame.Shares} condition, and the supplied point/numerical
construction are separate remaining obligations.
